% ECONOMICS_PROBLEM: {"id": "C78-213", "title": "Approximating the minimum repair budget for one-failure supermatches", "jel_code": "C78", "source_id": "OP-0560"}
% !TeX program = xelatex
\documentclass[11pt,a4paper]{article}
\usepackage[margin=23mm]{geometry}
\usepackage{fontspec}
\setmainfont{Latin Modern Roman}
\usepackage{amsmath,amssymb,mathtools}
\usepackage{xcolor,enumitem,booktabs,longtable,array,xurl,hyperref}
\definecolor{muted}{HTML}{53636C}
\hypersetup{colorlinks=true,urlcolor=blue,linkcolor=blue}
\setlength{\parindent}{0pt}
\setlength{\parskip}{5pt}
\newcommand{\R}{\mathbb R}
\newcommand{\E}{\mathbb E}
\newcommand{\N}{\mathbb N}
\newcommand{\ind}{\mathbf 1}
\newcommand{\Var}{\operatorname{Var}}
\newcommand{\Cov}{\operatorname{Cov}}
\newcommand{\supp}{\operatorname{supp}}
\DeclareMathOperator*{\argmin}{arg\,min}
\DeclareMathOperator*{\argmax}{arg\,max}
\newcommand{\entrymeta}[3]{{\sffamily\small\color{muted}\textbf{\texttt{#1}}\quad|\quad #2\par #3\par}}
\newcommand{\Problemref}[1]{\texttt{#1}}
\newcommand{\JELref}[2]{\texttt{#2} (JEL \texttt{#1})}
\begin{document}
\section*{Approximating the minimum repair budget for one-failure supermatches}
\textbf{Problem ID:} \texttt{C78-213}\quad \textbf{JEL:} \texttt{C78}\par
% BEGIN ECONOMICS STATEMENT
\label{problem:OP-0560}
\label{beyond:MM-19}
\entrymeta{MM-19}{Matching markets and allocation}{Published research direction; corrected setting}
\subsection*{Definitions and assumptions}
Let $I$ be a solvable strict stable-roommates instance: a finite set $A$ has mutually acceptable pairs, strict orders over acceptable partners, and every acceptable partner outranks unmatched. Stability excludes any acceptable pair whose endpoints both prefer one another to their assigned partners. Let $\mathcal S(I)$ be the stable matchings; fixed pairs belong to every one of them. For $M\in\mathcal S(I)$ let $B(M)$ be its non-fixed pairs and define
\[
 b_I(M)=\max_{e\in B(M)}\ 
 \min_{\substack{N\in\mathcal S(I)\\e\notin N}}
 \bigl(|M\setminus N|-1\bigr),\qquad
 b^*(I)=\min_{M\in\mathcal S(I)}b_I(M).
\]
Take the empty maximum to be zero. Non-fixedness ensures that every inner feasible set is nonempty. Only removed original pairs are counted, and the failed pair is subtracted.
\subsection*{Formal question}
Determine the best polynomial-time approximation guarantees for minimizing $b_I(M)$ over stable matchings. Specifically, seek an output $M$ satisfying $b_I(M)\le\rho b^*(I)$, with an explicit zero-optimum convention, or a stated additive bound. Classify the corresponding optimization problem in strict stable marriage as well, distinguishing that bipartite restriction from roommates.
\subsection*{Scope and status}
The source's immediate approximation research direction concerns stable roommates, where its cited work offers heuristics and polynomial one-failure verification. The supplied catalog's unqualified notation $I$ has therefore been narrowed to this primary setting, with marriage stated as a related restriction. The objective is equivalent to minimizing the budget $b$ for which a $(1,b)$-supermatch exists under the displayed pair convention. Instances admitting no stable matching are outside the optimization domain and should be reported as infeasible, rather than assigned an unexplained numerical optimum. The theoretical approximation status is attributed to the 2019 source.
\subsection*{References for the source of the problem}
Katar\'ina Cechl\'arov\'a, \'Agnes Cseh, and David F. Manlove (2019), \emph{Selected open problems in Matching Under Preferences}, Section~3.3.2, paragraph on metaheuristics and unstudied approximability: \url{https://hpi.de/friedrich/docs/publications/2019/BEATCS.pdf}. Begum Genc, Mohamed Siala, Gilles Simonin, and Barry O'Sullivan (2019), \emph{An Approach to Robustness in the Stable Roommates Problem and Its Comparison with the Stable Marriage Problem}, CPAIOR 2019, pp.~320--336, Definition~1 and optimization discussion: \url{https://siala.github.io/preprints/19-CPAIOR.pdf}.

\subsection*{Common conventions: Source mathematical conventions}

\textbf{Well-defined objects.}
All probability spaces and measurable variables are specified or inherited
from a local statistical model. Expectations used as finite targets require
the stated integrability. A decision rule or estimator is measurable with
respect to its available information; randomized rules may use an auxiliary
independent random variable. Norms without a subscript are Euclidean in
finite-dimensional spaces. An optimizer is requested only when attainment
is assured; otherwise the question concerns the optimal value and
$\varepsilon$-optimal admissible rules for every $\varepsilon>0$.

\textbf{Identification.}
A structural model $m\in\mathcal M$ implies an observable law
$P_m^{\rm obs}$ and target $\theta(m)$. The target is point identified on
$\mathcal M$ if
\[
 P_{m_1}^{\rm obs}=P_{m_2}^{\rm obs}
 \quad\Longrightarrow\quad\theta(m_1)=\theta(m_2)
 \qquad(m_1,m_2\in\mathcal M).
\]
When this fails, the sharp identified set is
\[
 \Theta(P;\mathcal M)=\{\theta(m):m\in\mathcal M,
                                  P_m^{\rm obs}=P\}.
\]
Specifying an intervention or estimand does not establish identification.
Positivity, exclusion, causal sufficiency, measurement restrictions, and
transport assumptions are substantive local inputs, never automatic
consequences of notation.

\textbf{Minimax risk and nontrivial inference.}
For a statistical experiment with data law $P^{(n)}$, target $\theta(P)$,
and nonnegative loss $\ell$, define
\[
 \mathcal R_n(\mathcal P,\ell)
   =\inf_{\widehat\theta_n}\sup_{P\in\mathcal P}
      \E_{P^{(n)}}\ell(\widehat\theta_n,\theta(P)).
\]
The infimum is over measurable estimators. A sharp rate is a sequence
$r_n$ for which $0<c\leq C<\infty$, independent of $n$, satisfies
$c r_n\leq\mathcal R_n\leq C r_n$ for all sufficiently large $n$.
Squared-error parametric risk is of order $n^{-1}$; the corresponding
root-mean-squared error is of order $n^{-1/2}$. These different conventions
are distinguished locally. A uniform asymptotic confidence region $C_n$
must satisfy
\[
 \liminf_{n\to\infty}\inf_{P\in\mathcal P}
     \Pr_{P^{(n)}}\{\theta(P)\in C_n\}\geq 1-\alpha,
 \qquad 0<\alpha<1,
\]
together with an informative diameter or expected-width requirement.
Coverage alone permits the vacuous whole parameter space. Testing questions
similarly impose both size control and power against a declared separated
alternative class.

\textbf{Binary causal baseline.}
When an entry explicitly invokes this baseline, one observes independent
copies of $O=(X,W,Y)$, with $W\in\{0,1\}$ and potential outcomes $Y(0),Y(1)$.
Assume consistency $Y=Y(W)$, no interference, conditional exchangeability
$(Y(0),Y(1))\perp W\mid X$, and overlap
$\epsilon\leq e(X)\leq1-\epsilon$ almost surely for a fixed
$0<\epsilon<1/2$, where
\[
 e(x)=\Pr(W=1\mid X=x),\qquad
 m_w(x)=\E[Y\mid W=w,X=x],\qquad
 \tau(x)=m_1(x)-m_0(x).
\]
These assumptions identify conditional mean effects almost everywhere
under the covariate law. Pointwise or uniform effects require appropriate
regularity and a specified support. Continuous treatment, longitudinal
data, adaptive experiments, and network interference require their own
assumptions in place of this baseline. Bounded outcomes, densities, and
smoothness are additional assumptions only where stated.

\textbf{Function classes.}
On $[0,1]^d$, $\mathcal H_d^s(L)$ is the isotropic Hölder ball of radius
$L>0$: put $k=\lceil s\rceil-1$ and $\gamma=s-k\in(0,1]$, bound derivatives
through order $k$, and require the order-$k$ derivatives to be
$\gamma$-Hölder with constant at most $L$. Thus integer orders use a
Lipschitz final derivative convention. The class
$\operatorname{BV}_d(L)$ consists of functions $f\in L^1((0,1)^d)$ whose
distributional derivative is a finite vector measure and for which
$\|f\|_{L^1}+|Df|((0,1)^d)\leq L$.
An $L^\infty$ bound or a particular pointwise version is imposed separately
when needed. Sparse and neural-network classes require a declared
dictionary or architecture and coefficient bounds.

An anisotropic H\"older class with declared block smoothness indices means
coordinatewise H\"older regularity in each block, uniformly in the other
blocks, with all corresponding derivative and H\"older bounds at most
the stated radius. Derivative orders and the integer-order convention in
each block are those just given. Mixed-derivative bounds are additional
restrictions only when explicitly stated.

\textbf{Decision and computational questions.}
Policy regret compares admissible feasible policies under the same law and
constraints. In computational entries, finite data and thresholds have
rational encodings; running time is measured in their total bit length.
Real-valued equilibrium search uses the explicitly declared algebraic or
FIXP encoding. An approximation ratio for maximization is written
$\operatorname{OPT}/\operatorname{ALG}$ and for minimization as
$\operatorname{ALG}/\operatorname{OPT}$, with zero optima and infeasible
instances handled locally. Historical approximation bounds are not
automatically current bounds.
% END ECONOMICS STATEMENT
\end{document}
